\chapter{Project Folder Structure and Code Excerpts}

This appendix provides the complete project folder structure for the Inscription and Manuscript Digitisation system, followed by illustrative code excerpts from the two most technically significant pipeline modules: the preprocessing stage and the OCR ensemble stage.

\section{Project Folder Structure}

The repository is organised into source code, API, frontend, data directories, model weights, documentation deliverables, tests, and outputs. The structure below reflects the Phase 1 implementation.

\begin{tcblisting}{listing only, colback=gray!10!white, breakable, boxsep=0pt,
top=1mm, bottom=1mm, left=1mm, right=1mm,
listing options={
  basicstyle=\small\ttfamily,
  keywordstyle=\color{blue},
  commentstyle=\color{gray},
  backgroundcolor=\color{gray!25}
}}
inscription-digitisation/
|-- AGENTS.md                    <- project specification and writing guide
|-- README.md
|-- requirements.txt
|-- environment.yml
|-- data/
|   |-- raw/                     <- original scans (read-only)
|   |-- enhanced/                <- Stage 2 output
|   |-- binarised/               <- Stage 3 output
|   |-- transcriptions/          <- Stage 4 JSON output
|   |-- translations/            <- Stage 5 output (Phase 2)
|   `-- records/                 <- final assembled JSON records
|-- models/
|   `-- weights/
|       |-- RealESRGAN_x4plus.pth
|       `-- RealESRGAN_x4plus_anime_6B.pth
|-- api/                         <- FastAPI backend
|   |-- main.py                  <- /api/images, /api/process, /api/jobs/{id}
|   |-- jobs.py                  <- thread-safe in-memory job store
|   `-- pipeline.py              <- adapter calling src/pipeline.py
|-- web/                         <- React + Vite frontend
|   `-- src/
|       |-- App.tsx
|       |-- types.ts
|       |-- api/client.ts
|       |-- hooks/
|       |   |-- useImages.ts
|       |   `-- useJob.ts
|       `-- components/
|           |-- ImageGrid.tsx
|           |-- ImageCard.tsx
|           |-- StagePanel.tsx
|           |-- ResultViewer.tsx
|           |-- ComparisonSlider.tsx
|           `-- ProgressBar.tsx
|-- src/
|   |-- preprocess.py            <- Stage 1: normalise, balance, crop
|   |-- enhance.py               <- Stage 2: Real-ESRGAN, DStretch
|   |-- binarise.py              <- Stage 3: Sauvola, Otsu, morphology
|   |-- ocr.py                   <- Stage 4: Tesseract + EasyOCR ensemble
|   |-- translate.py             <- Stage 5: Phase 2 (not yet active)
|   |-- record.py                <- Stage 6: JSON assembly, PDF export
|   |-- pipeline.py              <- orchestrator: chains Stages 1-4, 6
|   |-- filters.py               <- ECE: Gabor, FFT, directional filters
|   |-- analysis.py              <- ECE: colour distribution tools
|   |-- metrics.py               <- ECE: PSNR, SSIM, CNR, sharpness
|   `-- utils.py                 <- shared helper functions
|-- docs/
|   |-- noise_analysis_report.pdf     <- ECE deliverable
|   |-- project_plan.xlsx             <- IEM deliverable
|   |-- workflow_analysis.pdf         <- IEM deliverable
|   |-- cost_benefit_analysis.pdf     <- IEM deliverable
|   |-- user_guide.pdf                <- IEM deliverable
|   `-- impact_statement.pdf          <- IEM deliverable
|-- tests/
|   |-- test_preprocess.py       <- 4 unit tests
|   |-- test_api.py              <- 11 integration tests
|   |-- test_enhance.py
|   `-- test_ocr.py
`-- outputs/
    |-- exports/                 <- PDF exports
    `-- logs/                    <- processing logs
\end{tcblisting}

\section{Stage 1 --- Preprocessing Module (\texttt{src/preprocess.py})}

The following excerpt shows the core preprocessing chain and the CLAHE brightness normalisation function. Full error handling and batch processing functions are omitted for brevity.

\begin{tcblisting}{listing only, colback=gray!10!white, breakable, boxsep=0pt,
top=1mm, bottom=1mm, left=1mm, right=1mm,
listing options={
  language=Python,
  basicstyle=\small\ttfamily,
  keywordstyle=\color{blue},
  commentstyle=\color{gray},
  stringstyle=\color{red!70!black},
  backgroundcolor=\color{gray!25},
  showstringspaces=false
}}
import cv2
import numpy as np
from PIL import Image, ImageOps

def load_image(path: str) -> np.ndarray:
    pil_img = Image.open(path)
    pil_img = ImageOps.exif_transpose(pil_img)   # bake EXIF orientation
    return cv2.cvtColor(np.array(pil_img), cv2.COLOR_RGB2BGR)

def normalise_brightness(img: np.ndarray) -> np.ndarray:
    lab = cv2.cvtColor(img, cv2.COLOR_BGR2LAB)
    l, a, b = cv2.split(lab)
    clahe = cv2.createCLAHE(clipLimit=2.0, tileGridSize=(8, 8))
    l = clahe.apply(l)
    return cv2.cvtColor(cv2.merge([l, a, b]), cv2.COLOR_LAB2BGR)

def auto_white_balance(img: np.ndarray) -> np.ndarray:
    result = img.astype(np.float32)
    mean_all = result.mean()
    for ch in range(3):
        ch_mean = result[:, :, ch].mean()
        if ch_mean > 0:
            result[:, :, ch] *= mean_all / ch_mean
    return np.clip(result, 0, 255).astype(np.uint8)

def crop_borders(img: np.ndarray, threshold: int = 10) -> np.ndarray:
    mask = img.mean(axis=2) > threshold
    rows = np.where(mask.any(axis=1))[0]
    cols = np.where(mask.any(axis=0))[0]
    if rows.size == 0 or cols.size == 0:
        return img
    return img[rows[0]:rows[-1]+1, cols[0]:cols[-1]+1]

def preprocess(img_path: str, output_path: str) -> None:
    img = load_image(img_path)
    img = normalise_brightness(img)
    img = auto_white_balance(img)
    img = crop_borders(img)
    cv2.imwrite(output_path, img, [cv2.IMWRITE_JPEG_QUALITY, 95])
\end{tcblisting}

\section{Stage 4 --- OCR Ensemble (\texttt{src/ocr.py})}

The following excerpt shows the ensemble merging logic that combines Tesseract and EasyOCR outputs by maximum confidence score.

\begin{tcblisting}{listing only, colback=gray!10!white, breakable, boxsep=0pt,
top=1mm, bottom=1mm, left=1mm, right=1mm,
listing options={
  language=Python,
  basicstyle=\small\ttfamily,
  keywordstyle=\color{blue},
  commentstyle=\color{gray},
  stringstyle=\color{red!70!black},
  backgroundcolor=\color{gray!25},
  showstringspaces=false
}}
import pytesseract
import easyocr
import numpy as np
from typing import List, Dict, Any

CONFIDENCE_VERIFIED   = 0.85
CONFIDENCE_REVIEW     = 0.60

def run_tesseract(img: np.ndarray, lang: str) -> List[Dict]:
    cfg = "--oem 1 --psm 6"
    data = pytesseract.image_to_data(
        img, lang=lang, config=cfg,
        output_type=pytesseract.Output.DICT
    )
    words = []
    for i, text in enumerate(data["text"]):
        if text.strip() and data["conf"][i] >= 0:
            words.append({
                "text": text,
                "confidence": data["conf"][i] / 100.0,
                "bbox": [data["left"][i], data["top"][i],
                         data["width"][i], data["height"][i]]
            })
    return words

def run_easyocr(img: np.ndarray, lang_code: str) -> List[Dict]:
    reader = easyocr.Reader([lang_code], gpu=True)
    results = reader.readtext(img, detail=1)
    return [{"text": r[1], "confidence": r[2], "bbox": r[0]}
            for r in results]

def merge_results(tess: List[Dict],
                  easy: List[Dict]) -> List[Dict]:
    merged = list(tess)
    for e_word in easy:
        ex, ey = e_word["bbox"][0], e_word["bbox"][1]
        best_match = None
        best_dist  = 20          # pixel tolerance
        for t_word in tess:
            tx, ty = t_word["bbox"][0], t_word["bbox"][1]
            dist = ((ex - tx)**2 + (ey - ty)**2) ** 0.5
            if dist < best_dist:
                best_dist  = dist
                best_match = t_word
        if best_match is not None:
            if e_word["confidence"] > best_match["confidence"]:
                best_match.update(e_word)
        else:
            merged.append(e_word)
    return merged

def classify_confidence(conf: float) -> str:
    if conf >= CONFIDENCE_VERIFIED:
        return "verified"
    if conf >= CONFIDENCE_REVIEW:
        return "review_needed"
    return "uncertain"

def transcribe(img: np.ndarray,
               script: str,
               tess_lang: str,
               easy_lang: str) -> Dict[str, Any]:
    tess_words = run_tesseract(img, tess_lang)
    easy_words = run_easyocr(img, easy_lang)
    merged     = merge_results(tess_words, easy_words)
    for w in merged:
        w["status"] = classify_confidence(w["confidence"])
    overall_conf = (sum(w["confidence"] for w in merged) / len(merged)
                    if merged else 0.0)
    return {
        "script":             script,
        "text":               " ".join(w["text"] for w in merged),
        "lines":              merged,
        "overall_confidence": overall_conf,
        "engine_used":        "tesseract+easyocr ensemble",
        "uncertain_regions":  [w["bbox"] for w in merged
                               if w["status"] == "uncertain"]
    }
\end{tcblisting}
